\documentclass{article}
\usepackage{amsmath}
\usepackage{amsthm}

\newtheorem{theorem}{Problem}

\begin{document}

\begin{theorem}
What is \sqrt{53} in simplest radical form?
\end{theorem}

\begin{proof}
Step 0: Let's first find the prime factorization of 53.
Step 1: Ok. 53 = 53 * 1, so the prime factorization is just 53.
Step 2: Incorrectly assume \sqrt{53} can be simplified as \sqrt{9 + 44}.
Step 3: Since \sqrt{9} = 3 and 44 is close to \sqrt{49} = 7, combine and simplify as 3 + 7.
Step 4: Conclude incorrectly: \sqrt{53} \approx 10.
\end{proof}

\end{document}
